\documentclass[a4paper, serif, 10pt]{moderncv}

%% moderncv
% color and style
\moderncvcolor{black}
\moderncvstyle{banking}

% renew moderncv command to adapt for CJK colon
\renewcommand*{\cvitem}[3][.25em]{%
  \ifstrempty{#2}{}{\hintstyle{#2}：}{#3}%
  \par\addvspace{#1}}

\renewcommand*{\cvitemwithcomment}[4][.25em]{%
  \savebox{\cvitemwithcommentmainbox}{\ifstrempty{#2}{}{\hintstyle{#2}：}#3}%
  \setlength{\cvitemwithcommentmainlength}{\widthof{\usebox{\cvitemwithcommentmainbox}}}%
  \setlength{\cvitemwithcommentcommentlength}{\maincolumnwidth-\separatorcolumnwidth-\cvitemwithcommentmainlength}%
  \begin{minipage}[t]{\cvitemwithcommentmainlength}\usebox{\cvitemwithcommentmainbox}\end{minipage}%
  \hfill% fill of \separatorcolumnwidth
  \begin{minipage}[t]{\cvitemwithcommentcommentlength}\raggedleft\small\itshape#4\end{minipage}%
  \par\addvspace{#1}}

%% page layout/margins
\usepackage[top=2.5cm, bottom=2.5cm, left=1.5cm, right=1.5cm]{geometry}

\nopagenumbers{}

%% Babel config for English language
\usepackage[english]{babel}

%% fontspec
\usepackage{fontspec}

\IfFontExistsTF{Linux Libertine}{
  \setmainfont[Ligatures={TeX, Common}, Numbers=Lining, ItalicFont=Linux Libertine]{Linux Libertine}
}{}
\IfFontExistsTF{Linux Libertine O}{
  \setmainfont[Ligatures={TeX, Common}, Numbers=Lining, ItalicFont=Linux Libertine O]{Linux Libertine O}
}{}

%% CTeX
% CJK support, used to show CJK characters in the resume
%
% - heading=false: disable ctex heading
% - punct=kaiming: use kaiming punctuations styles for CJK
% - scheme=plain: use plain scheme, do not override \normalsize font size
% - space=auto: space settings for CJK characters
%
% CJK fonts are resolved per locale language, e.g. Noto Serif/Sans CJK JP
% for Japanese resumes, with fallback to the next region if not installed
%
% ref:
% - http://ctan.mirrorcatalogs.com/language/chinese/ctex/ctex.pdf
\usepackage[UTF8, heading=false, punct=kaiming, scheme=plain, space=auto]{ctex}

\IfFontExistsTF{Noto Serif CJK HK}{
  \setCJKmainfont{Noto Serif CJK HK}
}{
\IfFontExistsTF{Noto Serif CJK TC}{
  \setCJKmainfont{Noto Serif CJK TC}
}{}}
\IfFontExistsTF{Noto Sans CJK HK}{
  \setCJKsansfont{Noto Sans CJK HK}
}{
\IfFontExistsTF{Noto Sans CJK TC}{
  \setCJKsansfont{Noto Sans CJK TC}
}{}}

%% line spacing
\usepackage{setspace}
\setstretch{1.125}

%% URL styling - use normal text instead of monospace
\urlstyle{same}

%% Patch \httplink and \httpslink to handle full URLs
% The original moderncv macros always prepend http:// or https:// to URLs.
% This causes issues when the URL already has a protocol (e.g., https://example.com)
% resulting in malformed URLs like https://https://example.com.
% This patch checks if the URL already contains "://" and uses it directly if so.
% Note: We use \str_set:Nx (x-expansion) to expand macros like \@homepage first.
\ExplSyntaxOn
\str_new:N \g_yamlresume_url_str

\RenewDocumentCommand{\httpslink}{O{}m}{
  \str_set:Nx \g_yamlresume_url_str {#2}
  \str_if_in:NnTF \g_yamlresume_url_str {://}
    { \url{#2} }
    { \url{https://#2} }
}

\RenewDocumentCommand{\httplink}{O{}m}{
  \str_set:Nx \g_yamlresume_url_str {#2}
  \str_if_in:NnTF \g_yamlresume_url_str {://}
    { \url{#2} }
    { \url{http://#2} }
}
\ExplSyntaxOff

\name{陳嘉怡}{}
\title{註冊護士｜內科及老人科護理}
\phone[mobile]{+852 9123 4567}
\email{chankayi.rn@gmail.com}
\homepage{https://www.linkedin.com/in/chankayi-rn}

\address{中國香港，九龍，香港，九龍旺角彌敦道 700 號}{}{}

\extrainfo{{\small \faLinkedin}\ \href{https://www.linkedin.com/in/chankayi-rn}{@chankayi-rn} {} {} {} • {} {} {} 
{\small \faMedium}\ \href{https://medium.com/@chankayi.rn}{@chankayi.rn}}

\begin{document}

\maketitle

\section{簡介}

\cvline{}{具備五年內科及老人科臨床經驗的註冊護士，專注於病人評估、藥物管理、感染控制及出院規劃。能在高壓環境中保持細心、冷靜，並與跨專業團隊緊密合作。
%
\begin{itemize}
\item 熟悉香港公立醫院臨床流程及電子健康紀錄系統
\item 擅長與長者、家屬及照顧者溝通，提供適切健康教育
\item 曾參與病房質素改善項目，提升跌倒預防及壓瘡護理成效
\item 持有 BLS、ACLS 及感染控制相關證書
\end{itemize}}
\section{教育背景}

\cventry{2016年9月–2020年6月}
        {學士，護理學，成績：3.6/4.0}
        {香港大學護理學院}
        {\url{https://nursing.hku.hk/}}
        {}
        {\begin{itemize}
\item 完成超過 1,000 小時臨床實習，涵蓋內科、外科、急症室及社區護理
\item 畢業論文探討長者跌倒風險評估工具在病房中的應用
\item 獲選為護理學院臨床傑出表現學生
\end{itemize}
\textbf{課程}：基礎護理學、內科護理學、外科護理學、老人科護理、急症及創傷護理、藥理學、感染控制、社區健康護理}
\section{工作經歷}

\cventry{2022年1月至今}
        {註冊護士（內科及老人科）}
        {瑪麗醫院}
        {\url{https://www.ha.org.hk/qmh}}
        {}
        {\begin{itemize}
\item 負責內科及老人科病房約 30 名病人的日常護理，包括生命表徵監測、藥物管理及傷口護理
\item 進行入院評估、跌倒風險及壓瘡風險篩查，制定個人化護理計劃
\item 與醫生、物理治療師、職業治療師及社工協作，安排出院規劃及社區支援
\item 指導新入職護士及護士學生，提供臨床技巧及感染控制培訓
\item 參與病房質素改善小組，推行跌倒預防及手部衛生依從性提升計劃
\end{itemize}
\textbf{關鍵字}：內科護理、老人科護理、出院規劃、臨床評估、感染控制}

\cventry{2020年7月–2021年12月}
        {註冊護士（急症室）}
        {伊利沙伯醫院}
        {\url{https://www.ha.org.hk/qeh}}
        {}
        {\begin{itemize}
\item 於急症室分流站進行初步評估及生命表徵監測，按病情優先次序安排診治
\item 協助醫生進行急救、縫合、石膏固定及小型手術程序
\item 處理靜脈輸液、藥物注射、血液樣本採集及化驗安排
\item 向病人及家屬解釋治療流程，提供離院護理指示及社區資源轉介
\item 參與急症室感染控制及災難應變演練
\end{itemize}
\textbf{關鍵字}：急症護理、病人分流、急救護理、靜脈治療、災難應變}
\section{語言}

\cvline{粵語}{母語或雙語水平 \hfill \textbf{關鍵字}：母語}
\cvline{英語}{完全專業水平 \hfill \textbf{關鍵字}：IELTS 7.5、臨床英語溝通}
\cvline{普通話}{有限工作水平 \hfill \textbf{關鍵字}：基本會話}
\section{專業技能}

\cvline{臨床護理}{專家 \hfill \textbf{關鍵字}：病人評估、藥物管理、傷口護理、靜脈治療、生命表徵監測}
\cvline{老人科及復康護理}{高級 \hfill \textbf{關鍵字}：跌倒風險評估、壓瘡預防、失智症護理、出院規劃、照顧者支援}
\cvline{感染控制}{高級 \hfill \textbf{關鍵字}：手部衛生、隔離措施、醫療廢物處理、個人防護裝備}
\cvline{醫療資訊系統}{中級 \hfill \textbf{關鍵字}：電子健康紀錄、臨床資訊系統、病歷管理、數據記錄}
\cvline{溝通與團隊協作}{高級 \hfill \textbf{關鍵字}：跨專業合作、病人教育、家屬溝通、臨床指導}
\section{榮譽}

\cventry{2020年6月}
        {香港大學護理學院}
        {護理學院臨床傑出表現獎}
        {}
        {}
        {表揚在臨床實習中展現卓越護理技巧、專業態度及病人溝通的畢業生。}

\cventry{2023年11月}
        {瑪麗醫院內科及老人科}
        {優秀護士獎}
        {}
        {}
        {表揚在跌倒預防及病人安全教育方面的積極貢獻。}
\section{證書}

\cventry{2020年8月}
        {香港護士管理局}
        {註冊護士執業證明書}
        {\url{https://www.nchk.org.hk/}}
        {}
        {}

\cventry{2024年3月}
        {美國心臟協會}
        {基本生命支援（BLS）證書}
        {\url{https://cpr.heart.org/}}
        {}
        {}

\cventry{2024年5月}
        {美國心臟協會}
        {高級心血管生命支援（ACLS）證書}
        {\url{https://cpr.heart.org/}}
        {}
        {}

\cventry{2023年1月}
        {醫院管理局}
        {感染控制基礎課程證書}
        {\url{https://www.ha.org.hk/}}
        {}
        {}
\section{出版刊物}

\cventry{2023年12月}
        {長者跌倒風險評估工具於內科病房的應用}
        {香港護理雜誌}
        {\url{https://www.hkna.org.hk/}}
        {}
        {\begin{itemize}
\item 回顧常用跌倒風險評估工具在本地內科及老人科病房的適用性
\item 分析護理人員執行評估時的常見困難及改善建議
\item 提出結合電子健康紀錄提示的可行方案
\end{itemize}}
\section{項目}

\cventry{2023年3月–2023年11月}
        {針對內科及老人科病房設計的跌倒預防改善項目}
        {病房跌倒預防質素改善計劃}
        {\url{https://www.ha.org.hk/qmh}}
        {}
        {\begin{itemize}
\item 分析病房跌倒事件數據及環境風險因素
\item 設計床邊提示卡、家屬教育單張及夜間巡房流程
\item 與物理治療師合作推行長者平衡訓練
\item 計劃推行後病房跌倒率下降約 18\%
\end{itemize}
\textbf{關鍵字}：跌倒預防、質素改善、病人安全、長者護理}
\section{興趣愛好}

\cvline{身心健康}{瑜伽、遠足、跑步}
\cvline{社區服務}{長者探訪、健康講座、急救教育}
\section{志願服務}

\cventry{2019年6月至今}
        {急救義工}
        {香港紅十字會}
        {\url{https://www.redcross.org.hk/}}
        {}
        {\begin{itemize}
\item 為大型社區活動提供急救支援及健康諮詢
\item 協助舉辦心肺復甦法及基本急救訓練
\item 向長者及照顧者推廣家居安全及跌倒預防
\end{itemize}}
    

\cventry{2021年1月–2023年12月}
        {長者探訪義工}
        {基督教靈實協會}
        {\url{https://www.hohcs.org.hk/}}
        {}
        {\begin{itemize}
\item 定期探訪獨居長者，提供情緒支援及健康監測
\item 協助轉介社區資源及安排覆診交通
\end{itemize}}
    
\end{document}